\section{Result, conventions, and provenance}
\label{sec:result}

\begin{theorem}[Exact 238-state consistency machine]
\label{thm:main}
Let $M_{238}$ be the machine specified by
\file{machines/zfc_combined238.tm}, whose SHA-256 digest is
\begin{center}\small\ttfamily
909f29c493e0ccc5d6a481fd7a20a9a82\\
eb934ab2aa7df17b5ec3789255a8c08.
\end{center}
It has 238 working states, one binary tape indexed by $\mathbb Z$, and
only left and right moves. The head starts at position zero in state zero,
and every tape cell initially contains zero. State 238 is a separate
halting outcome and has no outgoing transitions. For the conventional
finite first-order ZFC proof predicate,
\begin{align}
 \ZFC&\vdash
 \Halt(M_{238},\blank)\longleftrightarrow
 \exists p\,\Proof_{\ZFC}(p,\bot), \label{eq:mainhalt}\\
 \ZFC&\vdash
 \neg\Halt(M_{238},\blank)\longleftrightarrow\Con(\ZFC).
 \label{eq:maincon}
\end{align}
The proof imports the established consistency calibration of the exact
ProveIt register source, identified in Section~\ref{sec:consistency}.
No assumption that ZFC is consistent is used in proving either
biconditional.
\end{theorem}

The user's threshold is fewer than 239 working states. This construction
meets it. The initializer is part of the 238 states; there is no preloaded
program, free input word, extra tape, extra alphabet symbol, or stationary
move. The distinction between working states and the halt outcome is the
same convention used by the supplied 254-state and ProveIt machines.

\begin{center}
\begin{tabular}{@{}L{0.42\textwidth}L{0.48\textwidth}@{}}
\toprule
Property & Exact value\\
\midrule
Working states & $0,\ldots,237$\\
Transition rows & $476$, two for each working state\\
Tape alphabet; blank symbol & $\{0,1\}$; $0$\\
Tape coordinates & $\mathbb Z$\\
Start state; initial head & $0$; $0$\\
Dispatcher; halt outcome & $84$; $238$\\
Inherited labels & Start $95$, dispatcher $96$, halt $254$\\
Exact initialization duration & $409,835$ transitions\\
Logical register seed & $(3,0,0,0,3,0,0,0,0)$\\
Source leaves; guarded interfaces & $6,455$; $8,386$\\
\bottomrule
\end{tabular}
\end{center}

A row \term{q b w D n} means: in state $q$, reading $b$, write $w$,
move one cell in direction $D\in\{L,R\}$, and enter state $n$.
The file's \term{\#! start 0} line fixes the initial state.
Appendix~\ref{sec:table} prints every transition.

\subsection{The two proof routes}

The primary route verifies the final table directly. It consists of exact
source identification, a source control-flow certificate, all unbounded
instruction interfaces for the new table, and an exact complete startup
frame. Composing these facts with the inherited source calibration proves
Theorem~\ref{thm:main}. It does not assume that the predecessor's machine
table was correctly compiled.

The second route proves a local reduction from the incoming Boundary239
machine. It checks the complete changed table, its incoming-state
invariant, four symbolic local paths, the full predecessor startup, and
the required guard over Boundary239's own unbounded runtime frames.
This route explains why the reductions compose and supplies a shorter
target for later formal verification.

Both routes use ordinary induction over naturals or finite lists. Python
calculations certify finite premises and symbolic identities; Python is
not a replacement for a proof-assistant kernel. The exact boundary is:
the mathematical construction and proofs are supplied, while a newly
compiled end-to-end Lean or Rocq theorem for this table is not supplied.

\subsection{Current repository and prior inputs}

The inspected ProveIt checkout is commit
\begin{center}\small\ttfamily
ce192939372ac7944fe2a0208c68a96d49fb5e3b.
\end{center}
The incoming commit \term{83fd3585484e5dffb37992f91f7dd4d3e70c437a}
contains three reports labeled 239 states. Their literal contents matter:
\begin{enumerate}
\item \file{ZFC_239_State_Construction_2026-10-19.zip} supplies
the distinct \emph{Boundary239} machine. It reparenthesizes the calibrated
source and reorganizes part of its decoder, eventually removing old
state 218.
\item \file{BusyBeaver239_Research_2026-10-19.zip} supplies the
\emph{coordinated-reuse} or \emph{delayed-erasure} machine, removing old
state 169.
\item \file{ZFC239_Companion_2026-10-19.zip} contains the
\emph{same literal delayed-erasure table} as the second archive.
These two reports contribute one state saving.
\end{enumerate}

All three are pinned in \file{PROVENANCE.json}, including archive digests.
The historical source pin used inside them is
\term{13ec39ed43148b1807f51ed6e6528e1c84659562}. The current
240-state table and its source data agree with those pinned inputs.
The present repository also reports an external Lean endpoint for the
exact 240-state binary machine; earlier statements in the incoming
reports about which external endpoint was then smallest are historical
and are not repeated as current facts \cite{rotation240,calibration}.

The supplied upstream work credits Mira-Cult for the 254-state machine
and Jeff Ketchersid for its Rocq verification. Its proof decomposes into
binary/register simulation, source enumeration, logical calculus, and
the native set theory's comparison with ZF/ZFC \cite{upstream}.
The Busy Beaver overview lists this 254-state milestone
\cite{bbwiki}. The present work improves the concrete state count using
the ProveIt line of reductions; it does not reassign credit for the
logical theory or claim a survey-based world-record certification.

\section{The concrete composition}
\label{sec:construction}

Use inherited labels throughout this section. The exact Boundary239
input is \file{reference/boundary239_old.tm}, with digest
\begin{center}\small\ttfamily
d061770f487937f9da652919c9bc9bb79\\
bc9bf137ca41dd7a08eb5ecda8e7719.
\end{center}
This historical file has no start directive; its documented inherited
start is 95. The new inherited-label file explicitly supplies that
directive.

Delete the two rows of state 169 and make precisely these three
retained-row edits:
\begin{center}
\begin{tabular}{@{}rll@{}}
\toprule
Row & Boundary239 & Combined238\\
\midrule
$80/0$ & $0R25$ & $0L106$\\
$91/1$ & $0L50$ & $1L167$\\
$147/0$ & $0R169$ & $0R80$\\
\bottomrule
\end{tabular}
\end{center}
All 473 other retained rows remain unchanged. The only incoming edge
to state 169 was $147/0$. Consequently no transition targets the
deleted state. There are now exactly 476 rows and 238 working states.
The canonical renumbering lists start 95 first, then every remaining
working label in increasing order, and sends halt 254 to 238.
The result is a bijection preserving every write, move, and target.

\begin{proposition}[Literal construction]
\label{prop:literal}
The three edits and one deletion above produce exactly the supplied
inherited-label table. Its canonical renumbering is exactly the table
in Theorem~\ref{thm:main}. Both presentations are total on their
working-state/symbol pairs, and their only nonworking target is the
designated halt outcome.
\end{proposition}
\begin{proof}
The constructor parses all predecessor rows, rejects duplicate or
malformed entries, checks the three old row values before editing, and
deletes state 169. It checks the complete resulting table against the
published bytes and both digests, then verifies every transition under
the saved state bijection. Its graph traversal also reaches every
working state and the halt outcome from the start in the directed
transition graph. This last fact is a graph property; it does not say
that the unique blank execution visits every state.
The finite checks are in \file{verification/check_construction.py}.
\end{proof}

\subsection{Why compatibility needed a new proof}

The boundary construction modifies 40 retained rows of the 240-state
table and deletes state 218. None of those edited rows occurs in the
three-edit list above. The helper transitions used by the delayed-erasure
paths also remain present. This makes the combination plausible, but
syntactic disjointness alone does not prove it correct.

The boundary transformation changes which paths the decoder follows.
A later shared helper could therefore be entered with a different
neighboring tape bit, invalidating a fold whose edited rows were
untouched. We prove the required invariant on the \emph{actual}
Boundary239 table and certify its additional right-neighbor guard
against the \emph{changed} source paths. The final table's direct
unbounded runtime proof supplies a second acceptance check.

The state saving comes from retaining a one temporarily at $91/1$.
Existing controls 167 and 151 then distinguish two caller contexts
using already present tape bits. Both continuations eventually erase
the intended cell and synchronize with the predecessor computation.
The method trades a few extra transitions for one working state.
No new data symbol or uncharged control is introduced.

\subsection{Exact state accounting}

\begin{center}
\begin{tabular}{@{}L{0.43\textwidth}rL{0.38\textwidth}@{}}
\toprule
Construction & States & Relevant change\\
\midrule
Upstream Mira-Cult table & 254 & Original verified benchmark\\
ProveIt Rotation240 & 240 & Calibrated source and binary baseline\\
Incoming Boundary239 & 239 & Source/decoder change; removes 218\\
Incoming delayed-erasure239 & 239 & Original source; removes 169\\
This Combined238 & \textbf{238} & Boundary source and both removals\\
\bottomrule
\end{tabular}
\end{center}

Source DAG counts are different quantities. The reparenthesization
changes 225 distinct source subtrees to 224 while retaining all 6,455
instruction occurrences. Its repaired binary implementation and fold
have their own charged state costs. The table above counts literal
working states, not source nodes, compressed bits, or a hypothetical
compiler output.
